\section{Forelesning 5: \emph{Digital innovasjon, disrupsjon og strategier}}
\label{sec:lec05}
\settermlabel{terms:lec05}

\subsection*{Oversikt}
Forelesning 5 (Sebastian Winslow) tar for seg \emph{digital innovasjon} som et
eget fenomen, distinkt fra alminnelig innovasjon, og kobler dette til
Christensens disrupsjonsteori. Hovedinnsikten: digitale komponenter har
egenskaper (homogenitet, redigerbarhet, selvreferering) som muliggjør stadig
\emph{rekombinasjon} -- og dermed nye former for innovasjon og disrupsjon.

\subsection{Hovedteori}

\subsubsection*{Definisjoner av innovasjon}
Pensum gir flere komplementære definisjoner:
\begin{itemize}
  \item \textbf{Rogers (1962/1983)}: \guillemotleft{}En idé, praksis eller gjenstand som
    oppfattes som ny av et individ eller en adopsjonsenhet.\guillemotright{}
  \item \textbf{Hukal \& Henfridsson (2017)}: \guillemotleft{}Samskapelse av nye tilbud
    gjennom \emph{rekombinasjon} av digitale og/eller fysiske komponenter.\guillemotright{}
  \item \textbf{Yoo et al.\ (2010)}: \guillemotleft{}Gjennomføring av nye kombinasjoner av
    digitale og fysiske komponenter for å produsere nye produkter.\guillemotright{}
  \item \textbf{Nambisan et al.\ (2017)}: \guillemotleft{}Bruk av digital teknologi under
    innovasjonsprosessen.\guillemotright{}
\end{itemize}

\begin{definition}{Digital innovasjon}
Gjennomføring av nye \emph{kombinasjoner} av digitale og fysiske komponenter
for å skape nye produkter, tjenester eller prosesser (Yoo et al.\ 2010).
\end{definition}

\subsubsection*{Tre egenskaper ved digital teknologi}
Det som gjør digital teknologi spesielt innovativ:
\begin{itemize}
  \item \term{Data homogeneity} -- all digital informasjon er bits; ulike
    datatyper kan kombineres uten kompatibilitetsproblemer.
  \item \term{Editability} -- digitalt innhold kan endres, også \emph{etter} det
    er distribuert.
  \item \term{Self-referentiality} -- digitale systemer kan referere til og
    endre seg selv; programvare kan modifisere programvare.
\end{itemize}

\subsubsection*{Rekombinasjon}
Digital innovasjon = kombinere eksisterende digitale + fysiske komponenter på
nye måter. \keypoint{Innovasjon trenger ikke være \guillemotleft{}nytt fra grunnen av\guillemotright{}; det
er ofte en \emph{ny kombinasjon} av eksisterende byggeklosser.}

\subsubsection*{Innovasjonstypologier}
Litteraturen skiller mellom flere innovasjonstyper som lett blandes sammen:
\begin{itemize}
  \item \term{Sustaining} -- forbedrer eksisterende verdiforslag for
    eksisterende kunder.
  \item \term{Radical} -- teknologisk gjennombrudd, stort sprang i ytelse.
  \item \term{Frugal} -- enkle, billige løsninger for ressursbegrensede
    markeder.
  \item \term{Disruptive} -- nedenfra-og-opp, etablerte aktører ignorerer den.
  \item \term{Discontinuous} -- bryter med tidligere paradigme/standarder.
  \item \term{Incremental} -- små, kontinuerlige forbedringer.
  \item \term{Architectural} -- omkonfigurering av eksisterende komponenter
    i ny systemarkitektur.
  \item \term{Business Model} -- innovasjon i hvordan verdi skapes, leveres
    og fanges.
\end{itemize}
Å være presis på hvilken type man snakker om er avgjørende -- en
\emph{disruptiv} innovasjon trenger ikke å være \emph{radikal} i teknologisk
forstand, og omvendt.

\begin{figure}[H]
\centering
\includegraphics[width=0.85\linewidth]{lec05-innovation-typologies.png}
\caption{Et kart over innovasjonstyper sentrert rundt \guillemotleft{}Innovation\guillemotright{}: Sustaining,
Radical, Frugal, Disruptive, Discontinuous, Incremental, Architectural,
Business Model.}
\end{figure}

\subsubsection*{Strategi-rammeverk for digital innovasjon (2$\times$2)}
Et 2$\times$2-rammeverk klassifiserer innovasjonstiltak langs aksene
\emph{Invent NEW vs.\ Enhance EXISTING} og \emph{INTERNAL vs.\ EXTERNAL
Innovation}. Hver kvadrant rommer ulike sub-temaer (Business Operations,
Employee Experience, Business Model, Customer Facing) -- nyttig for å
diagnostisere en virksomhets innovasjonsportefølje.

\begin{figure}[H]
\centering
\includegraphics[width=0.85\linewidth]{lec05-digital-innovation-strategy-framework.png}
\caption{2$\times$2-rammeverket: Invent New / Enhance Existing $\times$ Internal /
External Innovation.}
\end{figure}

\subsubsection*{Christensens disrupsjonsteori}

\begin{definition}{Disruptiv innovasjon}
En innovasjon som starter i bunnen av markedet eller i et forsømt segment
med et enklere, billigere tilbud som etablerte aktører ignorerer. Over tid
forbedrer disruptøren seg og beveger seg oppover, og fortrenger etablerte
aktører som har fokusert på sine beste kunder.
\end{definition}

Distinksjoner:
\begin{itemize}
  \item \term{Sustaining innovation} -- forbedrer eksisterende verdiforslag
    for eksisterende kunder.
  \item \term{Low-end disruption} -- enklere/billigere tilbud nedenfra.
  \item \term{New-market disruption} -- skaper et helt nytt marked av
    ikke-konsumenter.
\end{itemize}

\subsubsection*{Disrupsjons-\guillemotleft{}øye-modellen\guillemotright{}}
Slidene visualiserer disrupsjon som et 3$\times$3 (Type $\times$ Marked):
\begin{itemize}
  \item \emph{Type}: Disruptive Technology / Disruptive Radical / Disruptive
    Business Model Innovation.
  \item \emph{Marked}: Low-End / High-End / New Market Disruption.
\end{itemize}
Modellen krysset med Baiyere \& Salmela sin definisjon av disrupsjon gir et
presist språk for hvilken \emph{type} disrupsjon som er i spill.

\begin{figure}[H]
\centering
\includegraphics[width=0.85\linewidth]{lec05-disruption-eye-model.png}
\caption{\guillemotleft{}Øye-modellen\guillemotright{} -- Type $\times$ Marked-klassifisering av disrupsjon.}
\end{figure}

\subsubsection*{Digital innovation vs.\ digital disruption vs.\ disruptive innovation}
Tre lett forvekslede konsepter:
\begin{itemize}
  \item \term{Digital Innovation (DgI)} -- rekombinasjon av digitale +
    fysiske komponenter (Yoo, Hukal \& Henfridsson).
  \item \term{Digital Disruption (DgD)} -- disrupsjon \emph{forårsaket av}
    digital teknologi.
  \item \term{Disruptive Innovation (DI)} -- Christensens generiske
    disrupsjonsbegrep, ikke nødvendigvis digitalt.
\end{itemize}

\begin{figure}[H]
\centering
\includegraphics[width=0.7\linewidth]{lec05-digital-innovation-vs-disruption-venn.png}
\caption{Venn-diagram som disambiguerer DgI, DgD og DI. \emph{Mange digitale
innovasjoner er ikke disruptive; mange disrupsjoner er ikke digitale.}}
\end{figure}

\subsubsection*{Strategier mot/for digital disrupsjon}
Som i \autoref{sec:lec02}, en etablert aktør kan velge mellom fem grep:
\begin{itemize}
  \item \textbf{Forsvare} -- styrke eksisterende fortrinn og kunderelasjoner.
  \item \textbf{Investere} -- bygge nye digitale kanaler internt.
  \item \textbf{Kjøpe} -- oppkjøp av digitale aktører for å hente
    kapabiliteter.
  \item \textbf{Samarbeide} -- allianser og plattformpartnerskap.
  \item \textbf{Selv-disrupsjon} -- bygge en disruptiv enhet før noen andre
    gjør det.
\end{itemize}
Christensens råd til etablerte: opprett en separat enhet med eget mandat for å
forfølge disruptive muligheter -- den etablerte organisasjonen vil naturlig
avvise dem.

\subsection{Sentrale fagbegreper}\label{terms:lec05}
\begin{itemize}
  \item \term{Digital innovation (DgI)} -- rekombinasjon av digitale + fysiske
    komponenter.
  \item \term{Digital disruption (DgD)} -- disrupsjon forårsaket av digital
    teknologi.
  \item \term{Disruptive innovation (DI)} -- Christensens generiske
    disrupsjonsbegrep.
  \item \term{Recombination} -- nye kombinasjoner av eksisterende byggeklosser.
  \item \term{Data homogeneity}, \term{editability}, \term{self-referentiality}
    -- de tre egenskapene.
  \item \term{Sustaining innovation} -- forbedrer eksisterende posisjon.
  \item \term{Low-end} og \term{new-market disruption}.
\end{itemize}

\subsection{Diagnostiske spørsmål}
\begin{enumerate}
  \item Bruker initiativet eksisterende digitale + fysiske komponenter på en
    ny måte? $\rightarrow$ digital innovasjon.
  \item Forbedrer det det eksisterende verdiforslaget for eksisterende
    kunder? $\rightarrow$ sustaining.
  \item Treffer det en ikke-betjent eller forsømt kundegruppe med et
    enklere/billigere tilbud? $\rightarrow$ potensielt disruptivt.
  \item Er disrupsjonen \emph{teknologisk}, \emph{radikal} eller
    \emph{forretningsmodell-basert}? (øye-modellen)
\end{enumerate}

\subsection{Eksempler og caser}
\begin{eksempelboks}{Netflix vs.\ Blockbuster}
Klassisk disrupsjon: Netflix startet med et enklere postordretilbud, beveget
seg deretter oppover med streaming. Eksempel på \emph{new-market} som
utvidet seg oppover.
\end{eksempelboks}

\begin{eksempelboks}{Coop App: sustaining + DgI}
\emph{Sustaining innovation}, ikke disrupsjon. Forsterker det eksisterende
verdiforslaget (lojalitet, butikkopplevelse). Samtidig en \emph{digital
innovasjon}: rekombinasjon av lojalitetsdata (digital) + butikkbesøk (fysisk)
+ personlige tilbud (digital) + Scan \& Betal (digital + fysisk) $\rightarrow$
et nytt sammensatt tilbud.
\end{eksempelboks}

\begin{eksempelboks}{Potensielle disruptorer for Coop}
Rene netthandelsaktører som Nemlig (low-end på service); lavprisekspansjon
(Lidl, Netto -- low-end på pris); plattformaktører som tilbyr matplanlegging
og bestilling i stor skala (potensiell new-market disrupsjon).
\end{eksempelboks}

\subsection{Eksamensrelevans}
\begin{itemize}
  \item Definere digital innovasjon (Yoo et al.\ 2010) og forklare
    rekombinasjon.
  \item Forklare de tre egenskapene ved digital teknologi og hvorfor de
    muliggjør innovasjon.
  \item Forklare Christensens disrupsjonsteori.
  \item Disambiguere DgI / DgD / DI presist.
  \item Vurdere et konkret initiativ som sustaining vs.\ disruptiv.
  \item Knytte til \autoref{sec:lec04} (DT vs.\ ITEOT) -- en disruptiv
    innovasjon medfører ofte digital transformasjon; en sustaining
    innovasjon er ofte IT-enabled.
\end{itemize}
